% TOP_PROBLEM: {"id": 3086, "title": "Covering a square with unit squares", "category_id": 6, "category": {"id": 6, "name": "geometry", "display_name": "Geometry", "description": "Euclidean and non-Euclidean geometry, geometric structures.", "slug": "geometry", "order_index": 6, "created_at": "2026-07-31T15:26:25.671Z"}, "status": "open", "problem_number": "OPG-37327", "set_id": 12}
% FIRST_POSTED: 2026-10-01
% ATTEMPT_STATUS: unresolved
% COMPLETION: 5
% MODEL: GPT 6 Astra Ultra
% UnsolvedMath ID 3086; source catalogue code OPG-37327
% !TeX program = lualatex
% Research attempt, 2026-10-01. Canonical statement preserved verbatim.
\documentclass[11pt,a4paper]{article}
\usepackage[margin=25mm]{geometry}
\usepackage{fontspec}
\setmainfont{lmroman10-regular.otf}[BoldFont=lmroman10-bold.otf,
 ItalicFont=lmroman10-italic.otf,BoldItalicFont=lmroman10-bolditalic.otf]
\usepackage{amsmath,amssymb,mathtools}
\usepackage{unicode-math}
\setmathfont{latinmodern-math.otf}
\AtBeginDocument{\let\setminus\smallsetminus}
\usepackage{xurl,hyperref}
\hypersetup{colorlinks=true,urlcolor=blue}
\setcounter{secnumdepth}{0}
\setlength{\parindent}{0pt}
\setlength{\parskip}{5pt}
\setlength{\emergencystretch}{3em}
\begin{document}
\section*{Covering a square with unit squares}\label{3086}
{\small UnsolvedMath ID 3086 · OPG-37327 · Geometry · OpenGarden}
\subsection{Definitions and mathematical statement}
For $s>0$, let $Q_s=[0,s]^2\subset\mathbb R^2$.
A unit square is any set congruent to $[0,1]^2$ under a Euclidean isometry;
its sides need not be parallel to those of $Q_s$ or to those of the
other squares. A family $U_1,\ldots,U_m$ covers $Q_s$ when
$Q_s\subseteq\bigcup_{i=1}^m U_i$. Squares may overlap each other and
may extend beyond the boundary of $Q_s$.

The conjecture says that for every integer $n\ge1$ and every real
$s>n$, no family of $n^2+1$ unit squares covers $Q_s$. Equivalently,
\[
 Q_s\subseteq\bigcup_{i=1}^{n^2+1}U_i
       \quad\Longrightarrow\quad s\le n.
\]
Allowing fewer squares gives the same nonexistence assertion, since any
smaller cover can be augmented by unused or repeated squares. All boundaries
are included in the covering condition; covering almost every point by
area alone is not its definition.

The integer threshold is sharp at side length $n$ in the elementary
sense that an $n$ by $n$ array of unit squares covers $Q_n$ using only
$n^2$ squares. The question asks whether one additional square can enable
any positive increase in side length, however small. A pure area estimate
only yields $s^2\le n^2+1$, leaving a nonempty interval of possible side
lengths above $n$; controlling this interval requires the geometry of
covering rather than the sum of areas alone.
\subsection{Short English statement}
Can adding just one square to an $n$ by $n$ collection of unit squares ever cover a square with side strictly larger than $n$?
\subsection{Research attempt}
\paragraph{Scope and literature check.}
The original question [S2], checked 1 October 2026, allows rotations and
outward protrusion. This attempt proves the stronger counting obstruction
when all covering squares are axis-parallel, and quantifies precisely which
part of that obstruction survives when rotations are allowed. The geometric
certificate below does not rely on area or on assuming a tiling.

\paragraph{Separated-point certificate.}
Fix $s>n$ and put $\delta=s/n>1$. The set
\[
 T=\{(i\delta,j\delta):0\le i,j\le n\}\subseteq[0,s]^2
\]
has $(n+1)^2$ points. An axis-parallel unit square cannot contain two distinct
points of $T$: some coordinate of any such pair differs by at least
$\delta>1$, whereas each coordinate interval of the square has length one.
A cover must cover every point of $T$, including the boundary points.
It therefore needs at least $(n+1)^2>n^2+1$ squares. This proof allows arbitrary
overlaps, arbitrary translations, and arbitrary extension beyond the target.

\paragraph{A quantitative rotation obstruction.}
For a unit square rotated through angle $\theta$, reduce by its symmetries
to $0\le\theta\le\pi/4$. Both coordinate projections have length
$w(\theta)=\cos\theta+\sin\theta$. This follows by projecting its two edge
vectors: the width of their Minkowski sum is the sum of the projected widths.
If every covering square satisfies $w(\theta)<\delta$, the same point-count
argument still applies, using the squares' bounding boxes. Thus any cover
with only $n^2+1$ squares must have at least one square with
\[
 \cos\theta+\sin\theta\ge s/n.
\]
For $1<s/n\le\sqrt2$ this forces
\[
 \theta\ge\arcsin\!\left(\frac{s}{n\sqrt2}\right)-\frac\pi4.
\]
If $s/n>\sqrt2$, the projection test excludes the cover with no restriction
on angles. At equality the strict separation argument ceases to apply;
it is not legitimate to replace its strict inequalities by weak ones.
For small excess $s=n+\varepsilon$, the forced angle tends to zero with
$\varepsilon$, so this result provides no fixed positive angle barrier.

\paragraph{Failure of a rotation-invariant version.}
One might try to say that a unit square covers at most one of these grid
points regardless of orientation. That is false. A square rotated by
$\pi/4$ is a diamond whose horizontal diagonal has length $\sqrt2$.
It contains two points at horizontal separation $1+\varepsilon$ whenever
$0<\varepsilon\le\sqrt2-1$, after translation to their midpoint.
Thus the very local property needed by the discrete certificate fails in
the difficult interval near $s=n$. Replacing coordinate separation by
Euclidean separation greater than $\sqrt2$ restores the property but loses
the required $(n+1)^2$ points near this threshold.

\paragraph{Verification and first remaining gap.}
The obstruction counts exact points, not sampled uncovered area, so it is
unaffected by measure-zero cracks or boundary conventions. The projection
formula is checked at $\theta=0$ and $\pi/4$. The construction at $s=n$ by
$n^2$ squares is consistent: then grid separation is one and a single square
can contain several grid vertices. A general proof would need a certificate
that accounts simultaneously for the orientations of different squares,
not just a common bound on their coordinate widths. This attempt gives no
such certificate, nor a rotated cover contradicting the conjecture.

\paragraph{Final status.}
\textbf{UNRESOLVED.} The axis-parallel case and the stated bounded-projection
extension are proved; the unrestricted rotation case remains outside the
argument. No historical novelty is claimed. Completion estimate: 5\%.

\subsection{Sources}
\begin{itemize}
\item[S1] ULAM.AI and the UnsolvedMath Contributors, supplied problems.json, record 3086 (OPG-37327), OpenGarden collection. Catalogue identity and source transcription; CC BY 4.0.
\url{https://huggingface.co/datasets/ulamai/UnsolvedMath}.
\item[S2] Open Problem Garden, Covering a square with unit squares. Original problem and discussion; the mathematical formulation below is expanded from the supplied source record.
\url{http://www.openproblemgarden.org/op/covering_a_square_with_unit_squares}.

\end{itemize}
\end{document}
